% from ega3-1-fr.tex
\oldpage[0\textsubscript{III}]{14}
de $X$ admette un système fondamental de voisinages ouverts rétrocompacts
dans $X$ (et \emph{a fortiori} quasi-compacts)~; alors la condition pour une
partie $Z$ de $X$ d'être constructible dans $X$ est de nature \emph{locale},
autrement dit, il faut et il suffit que pour tout $x\in X$, il existe un
voisinage ouvert $V$ de $x$ tel que $V\cap Z$ soit constructible dans $V$. En
effet, si cette condition est vérifiée, il existe pour tout $x\in X$ un
voisinage ouvert $V$ de $x$ \emph{rétrocompact dans $X$} et tel que
$V\cap Z$ soit constructible dans $V$, en vertu de l'hypothèse sur $X$ et de
(\hyperref[0.9.1.8-fr]{9.1.8, (i)})~; il suffit alors de recouvrir $X$ par un
nombre fini de ces voisinages, et d'appliquer
(\hyperref[0.9.1.9-fr]{9.1.9}).
\end{env}

\begin{definition}[9.1.11]
\label{0.9.1.11-fr}
Soit $X$ un espace topologique. On dit qu'une partie $T$ de $X$ est
\emph{localement constructible dans $X$} si, pour tout $x\in X$, il existe
un voisinage ouvert $V$ de $x$ tel que $T\cap V$ soit constructible dans $V$.
\end{definition}

Il résulte aussitôt de (\hyperref[0.9.1.8-fr]{9.1.8, (i)}) que si $V$ est
tel que $V\cap T$ soit constructible dans $V$, alors pour tout ouvert
$W\subset V$, $W\cap T$ est constructible dans $W$. Si $T$ est localement
constructible dans $X$, alors, pour tout ouvert $U$ de $X$, $T\cap U$ est
localement constructible dans $U$, comme il résulte de la remarque
précédente. Cette même remarque montre que l'ensemble des parties localement
constructibles dans $X$ est stable par réunion finie et intersection finie~;
il est clair, d'autre part, qu'il est aussi stable par passage aux
complémentaires.

\begin{proposition}[9.1.12]
\label{0.9.1.12-fr}
Soit $X$ un espace topologique. Tout ensemble constructible dans $X$ est
localement constructible dans $X$. La réciproque est vraie si $X$ est
quasi-compact et si sa topologie admet une base formée d'ensembles
rétrocompacts dans $X$.
\end{proposition}

La première assertion résulte de la définition
(\hyperref[0.9.1.11-fr]{9.1.11}) et la seconde de
(\hyperref[0.9.1.10-fr]{9.1.10}).

\begin{corollary}[9.1.13]
\label{0.9.1.13-fr}
Soit $X$ un espace topologique dont la topologie admet une base formée
d'ensembles rétrocompacts dans $X$. Alors toute partie $T$ localement
constructible dans $X$ est rétrocompacte dans $X$.
\end{corollary}

En effet, soit $U$ un ensemble ouvert quasi-compact dans $X$~; $T\cap U$ est
localement constructible dans $U$, donc constructible dans $U$ en vertu de
(\hyperref[0.9.1.12-fr]{9.1.12}), et par suite quasi-compact en vertu de
(\hyperref[0.9.1.4-fr]{9.1.4}).

\subsection{Ensembles constructibles dans les espaces noethériens}
\label{subsection:0.9.2-fr}

\begin{env}[9.2.1]
\label{0.9.2.1-fr}
On a vu (\hyperref[0.9.1.7-fr]{9.1.7}) que dans un espace noethérien $X$, les
parties constructibles dans $X$ sont les \emph{réunions finies de parties
localement fermées de $X$}.

L'image réciproque d'un ensemble constructible dans $X$ par une application
continue d'un espace noethérien $X'$ dans $X$ est constructible dans $X'$.
Si $Y$ est une partie constructible d'un espace noethérien $X$, les parties
de $Y$ qui sont constructibles en tant que sous-espaces de $Y$ sont
identiques à celles qui sont constructibles en tant que sous-espaces de $X$.
\end{env}

\begin{proposition}[9.2.2]
\label{0.9.2.2-fr}
Soient $X$ un espace irréductible noethérien, $E$ une partie constructible de
$X$. Pour que $E$ soit partout dense dans $X$, il faut et il suffit que $E$
contienne une partie ouverte non vide de $X$.
\end{proposition}

